% Source: source/普通高考/2022/2022天津.pdf, p.1, question 4.
% Frequency-density histogram for six temperature intervals of width 0.20 C.
% Elements: arrowed axes; x ticks 13.55, ..., 14.75; y ticks 0,.40,.50,.60,.65,1.30,1.55;
% six contiguous white bars with heights .40,1.30,1.55,.60,.50,.65;
% dashed guides run from the y-axis to the left boundary of the corresponding bar.
% Relations: each bar encodes the source interval's frequency/group-width density;
% guides identify the labelled density levels, and the x-axis label gives global mean temperature.
\begin{tikzpicture}[baseline]
  \begin{axis}[exam statistical histogram,scaled ticks=false,
    width=16cm,
    height=15.429cm,
    xmin=-0.70, xmax=7.2,
    ymin=0, ymax=1.70,
    axis lines=middle,
    axis line style={-Stealth},
    xtick={0,1,2,3,4,5,6},
    xticklabels={13.55,13.75,13.95,14.15,14.35,14.55,14.75},
    ytick={0,0.40,0.50,0.60,0.65,1.30,1.55},
    % Place all y labels explicitly: the four low levels use alternating
    % columns/vertical offsets so 0.40--0.65 stay distinct at page scale.
    yticklabels={,},
    scaled y ticks=false,
    yticklabel style={/pgf/number format/fixed, /pgf/number format/precision=2,xshift=-4pt},
    tick align=outside,
    tick label style={font=\normalsize},
    yticklabel style={font=\normalsize, /pgf/number format/fixed, /pgf/number format/precision=2,xshift=-4pt},
    enlarge x limits=false,
    clip=false,
  ]
    % Bar interval values correspond to [13.55,13.75], ..., [14.55,14.75].
    \addplot[
      ybar interval,
      draw=black,
      fill=white,
      line width=0.55pt,
      area legend,
    ] coordinates {
      (0,0.40) (1,1.30) (2,1.55) (3,0.60)
      (4,0.50) (5,0.65) (6,0.65)
    };
    % Draw guides over the white bars, as in the source figure; each ends at
    % the left boundary of the bar carrying that density level.
    \addplot[black,dashed,line width=0.45pt] coordinates {(0,1.30) (1,1.30)};
    \addplot[black,dashed,line width=0.45pt] coordinates {(0,1.55) (2,1.55)};
    \addplot[black,dashed,line width=0.45pt] coordinates {(0,0.60) (3,0.60)};
    \addplot[black,dashed,line width=0.45pt] coordinates {(0,0.50) (4,0.50)};
    \addplot[black,dashed,line width=0.45pt] coordinates {(0,0.65) (5,0.65)};
    \examstatylabel{$\dfrac{\text{频率}}{\text{组距}}$};
    \examstatxlabel{全球年平均气温/$^{\circ}$C};
    % Explicit labels keep the baseline 0 and all close levels visible.
    \node[anchor=south east,font=\normalsize,inner sep=0pt,xshift=-6pt,yshift=4pt] at (axis cs:-0.08,0) {0};
    \node[anchor=east,font=\normalsize,inner sep=0pt,xshift=-4pt] at (axis cs:-0.080000,0.40) {0.40};
    \node[anchor=east,font=\normalsize,inner sep=0pt] at ([yshift=-1.2pt]axis cs:-0.08,0.50) {0.50};
    \node[anchor=east,font=\normalsize,inner sep=0pt] at ([yshift=-4.5pt]axis cs:-0.08,0.60) {0.60};
    \node[anchor=east,font=\normalsize,inner sep=0pt] at ([yshift=4.5pt]axis cs:-0.08,0.65) {0.65};
    \node[anchor=east,font=\normalsize,inner sep=0pt,xshift=-4pt] at (axis cs:-0.080000,1.30) {1.30};
    \node[anchor=east,font=\normalsize,inner sep=0pt,xshift=-4pt] at (axis cs:-0.080000,1.55) {1.55};
    % Keep the first decimal tick visible at the axis crossing.
    \node[anchor=north east,font=\normalsize,inner sep=0pt,yshift=-4pt] at (axis cs:0,0.000000) {13.55};
  \end{axis}
\end{tikzpicture}
